\section*{Chapter \ref{ch:m1:delta-one-instruments} --- Delta-One Instruments}

\begin{solution}{exo:m1:delta-one-instruments:1}
The fund receives $50\,000\,000 \times 6\% = \$3\,000\,000$ and pays
$50\,000\,000 \times 4.6\% \times 91/360 = \$581\,389$: a net receipt of
\$2\,418\,611.
\end{solution}

\begin{solution}{exo:m1:delta-one-instruments:2}
Zero spread: the fair value, 6\,043.5. Fifty basis points: $6\,000 \times (1 +
(4.7\% - 1.3\%) \times 0.25) = 6\,051$. Each index point on a three-month
future is worth about 6.7 basis points of annual financing.
\end{solution}

\begin{solution}{exo:m1:delta-one-instruments:3}
Parity $42.10 \times 2 \times 1.17 = \$98.51$; premium $98.90/98.51 - 1 = 39$
basis points.
\end{solution}

\begin{solution}{exo:m1:delta-one-instruments:4}
Shares: 10 of trading plus 50 of tax, no funding. CFD: 16 of trading plus 250
a year. $60 = 16 + 250\,y$ gives $y = 0.176$ year: the CFD is cheaper for
holding periods under 64 days.
\end{solution}

\begin{solution}{exo:m1:delta-one-instruments:5}
$1.5\% \times 30\% = 45$ basis points a year; at 15\%, 22.5. Since section
871(m) treats the dividend-equivalent payment of a \omterm{def:m1:delta-one-instruments:deltaone}{delta-one} swap on a US
stock as a US-source dividend, the dealer must withhold on it as on the
dividend: a swap that pays 100\% is either not \omterm{def:m1:delta-one-instruments:deltaone}{delta-one}, or mispriced
elsewhere (a wider \omterm{def:m1:delta-one-instruments:swap}{funding spread}), or relies on a tax position the holder
should want to understand before it is challenged.
\end{solution}

\begin{solution}{exo:m1:delta-one-instruments:6}
Per year, with cash: shares cost 30 (dividends); the swap $40 + 10 = 50$; the
future $6 + 30 + 30 = 66$; the fund $7 + 30 = 37$. Fixed: shares 60, swap 6,
future 2, fund 6. Shares against swap: $60 + 30y = 6 + 50y$, $y = 2.7$ years.
Fund against future: $6 + 37y = 2 + 66y$, $y = 4/29$ year, fifty days.
\end{solution}

\begin{solution}{exo:m1:delta-one-instruments:7}
Shares 80, fund 27.75, future 18.5, swap 18.5, CFD 86 basis points;
annualised 320, 111, 74, 74 and 344. Future and swap tie. In practice: the
future needs no documentation and has no counterparty exposure beyond the
clearing house, but exists only on the index and in fixed sizes and
expiries; the swap can be on a custom basket and pays more of the dividend,
but needs an agreement, a credit line and a dealer willing to quote.
\end{solution}

\begin{solution}{exo:m1:delta-one-instruments:8}
CFD: $16 + 3 \times 250 + 3 \times 30 = 856$ basis points. Shares: $10 + 50 +
3 \times 30 = 150$. The page omitted the financing, charged every day on the
whole notional even to a client who has the cash, and, by quoting ``no
commission'', the spread in which the commission now lives. The claim is
true for holding periods under about two months.
\end{solution}

\begin{solution}{pb:m1:delta-one-instruments:1}
\textbf{1.} $5.00 \times 4 \times 1.25 = \$25.00$.
\textbf{2.} $6 + 150 + 0.05/25 = 6 + 150 + 20 = 176$ basis points.
\textbf{3.} $-(6 + 20) = -26$ basis points.
\textbf{4.} $+120$ basis points: inside the band, nothing pays.
\textbf{5.} Creating receipts costs 1.5\% and destroying them does not, so
supply responds to discounts at once and to premiums only beyond 1.76\%: the
receipt should trade at a small premium on average, whenever US demand
grows.
\textbf{6.} Buy 400\,000 ordinaries: \pounds2\,000\,000 $=$ \$2\,500\,000.
Deposit charge \$37\,500, issue fee \$5\,000, trading \$1\,500: \$44\,000. Sell
100\,000 ADRs: \$2\,550\,000. Profit \$6\,000, 24 basis points ($200 - 176$).
\textbf{7.} Buy ADRs \$2\,490\,000, cancel fee \$5\,000, trading \$1\,500, sell
the ordinaries for \$2\,500\,000: profit \$3\,500, 14 basis points ($40 - 26$).
\textbf{8.} $8\% \times \sqrt{2/252} = 71$ basis points, several times either
profit: the currency must be sold forward at the moment of the trade. The
conversion is an arbitrage only with three legs.
\textbf{9.} Both. Issuing: the firm sells ADRs it will receive in two days
and must borrow ADRs to deliver. Cancelling: it sells ordinaries it will
receive in two days and must borrow them in London. The borrow fee and
availability belong in the band.
\textbf{10.} $25.00 \times (1 + 1.1 \times 1.5\%) = \$25.41$.
\textbf{11.} 180 basis points to the stale parity; 15 to the estimate.
\textbf{12.} He sells the ADR and buys nothing: London is closed. He is
short the stock overnight against a hedge in US index futures at best; what
he has sold is 15 basis points of premium and the whole of the stock's
idiosyncratic risk until 08:00 London time.
\textbf{13.} Regress the ordinary's return from the London close to the next
London open on the ADR's return (currency-adjusted) from 11:30 to 16:00 New
York time, over a few years, controlling for the US index; the coefficient
should be close to one and the $R^2$ high.
\textbf{14.} Because everyone in the opening auction has seen the ADR: the
\omterm{def:m1:auctions:imbalance}{indicative price} already contains it. The predictability is of the open
relative to yesterday's close, not relative to any price at which one can
trade.
\textbf{15.} The ADR costs the depositary's custody fee, a few cents per
receipt a year, and (since it was created once) embeds the 1.5\% only if she
is the one creating it; it spares her a UK custody account, sterling
settlement, and currency conversion of each dividend. The ordinary costs
0.5\% on purchase and the foreign plumbing, and has the deeper market.
\textbf{16.} There is no conversion left, so no band: one share, one
register, with London trading becoming a secondary listing or disappearing.
\textbf{17.} A relative-value position with no mechanism of convergence: it
profits only if the restriction is lifted or if investors in the two markets
change their minds. It should be sized as a directional bet on a spread.
\textbf{18.} Dealers with an account at the depositary, custody in both
markets, intermediary status for the tax, stock borrow in both lines and a
currency desk: banks' \omterm{def:m1:delta-one-instruments:deltaone}{delta-one desks} and a few trading firms.
\textbf{19.} \textbf{From $-26$ to $+176$ basis points.}
\textbf{20.} The same exposure in the wrapper that is cheapest for the
client, priced from differences in funding, tax and access that the desk can
bear more cheaply than the client can.
\end{solution}

\begin{solution}{iq:m1:delta-one-instruments:1}
\omterm{def:m1:delta-one-instruments:deltaone}{Delta-one} instruments carry an exposure without optionality: shares, funds,
futures, swaps, receipts. They all give the same return before costs, so
the desk's business is the costs: it quotes swaps and funds clients'
positions, makes markets in funds and receipts, arbitrages futures against
baskets, and manages the dividend, borrow and financing risks left over.
\iqlookfor{``same exposure, different costs'', and at least three concrete
activities.}
\end{solution}

\begin{solution}{iq:m1:delta-one-instruments:2}
Cost the five components for her horizon. The fund: spread, \omterm{def:m1:the-buy-side:fees}{management fee},
withholding inside the fund. The future: tiny commissions, four rolls a
year, and the implied financing spread, which she pays even if she has the
cash. A funded long-term holder is usually better off in the fund; a
leveraged or short-term one in the future; and the answer moves with the
richness of the roll.
\iqlookfor{the implied financing spread as the future's hidden fee.}
\end{solution}

\begin{solution}{iq:m1:delta-one-instruments:3}
Long holders of futures pay 80 basis points a year over the benchmark;
anyone who can hold the shares and sell the future earns it. The trade is
cash-and-carry: buy the basket, sell the future, finance the basket. What
stops it is balance sheet: the position is large, low-yielding and consumes
dealer capital, most of all at year-end, which is exactly when the spread is
widest. A funded investor can also collect it passively, by replacing
futures with a fund.
\iqlookfor{the arbitrage \emph{and} the reason it is limited.}
\end{solution}

\begin{solution}{iq:m1:delta-one-instruments:4}
\omterm{def:m1:delta-one-instruments:swap}{Funding spread}: my own cost of financing the hedge, the capital and
balance-sheet charge, the client's credit and margin terms, any transaction
tax I cannot avoid, less what the shares earn me in stock lending.
Dividend percentage: my net-of-tax receipt on the dividend given where I
book the hedge, less a margin. For a hard-to-borrow stock on the short side,
the borrow fee. Competition sets how much of each benefit I pass on.
\iqlookfor{balance sheet, tax position and lending revenue, all three.}
\end{solution}

\begin{solution}{iq:m1:delta-one-instruments:5}
Only if conversion is possible and cheaper than 1\%. Check: is the home
market open (otherwise parity is stale); the ratio and the currency rate
used; issue fee per receipt; any deposit tax; whether the programme is open
for issuance (some are capped or closed); borrow in the receipt; settlement
timing and the currency hedge. If all pass, issue receipts against
ordinaries.
\iqlookfor{stale parity and the closed-programme case.}
\end{solution}

\begin{solution}{iq:m1:delta-one-instruments:6}
Financing becomes explicit: a spread on the long notional and on the short
notional, not an interest rate on net cash. No purchase tax, which can
matter more than everything else in taxed markets. Dividends at the swap's
percentage, on both sides. Short availability and fees as the dealer quotes
them, with recall replaced by the dealer's right to terminate or reprice.
Resets realise P\&L in cash monthly. And capacity is the dealer's balance
sheet, which can be withdrawn.
\iqlookfor{long and short spreads charged on gross notional.}
\end{solution}
